\documentclass[11pt]{article}
\usepackage[margin=0.78in]{geometry}
\usepackage{amsmath,amssymb,bm}
\usepackage{booktabs}
\usepackage{enumitem}
\usepackage{hyperref}
\hypersetup{colorlinks=true, linkcolor=black, urlcolor=black}
\setlength{\parindent}{0pt}
\setlength{\parskip}{0.55em}

\title{Cyrus Guided Notes: Stanford CS229 Machine Learning}
\author{Stanford course pack}
\date{2026-05-15}

\begin{document}
\maketitle

\section*{Course source}
\textbf{Official source:} \url{https://cs229.stanford.edu/}

\textbf{Why this course is in the system.} Empirical risk, gradient descent, generative models, SVMs, neural networks, and generalization.

\textbf{Learning output.} Next batch target for a machine-learning formula and problem route.

\section*{Prerequisite checkpoint}
\begin{itemize}[leftmargin=1.4em]
  \item Calculus
  \item Linear algebra
  \item Probability
\end{itemize}

\section*{Formula checkpoint}
Do not memorize these first. Read each formula as a sentence: what changes, what is observed, what is optimized, and what output it produces.
\begin{align*}
  \min_{\theta}\frac{1}{N}\sum_{i=1}^{N}\ell(f_{\theta}(\bm{x}_i),y_i)+\lambda\Omega(\theta)\\
  \theta\leftarrow\theta-\alpha\nabla_{\theta}J(\theta)\\
  p(y=1|\bm{x})=\sigma(\theta^{\mathsf T}\bm{x})\\
  p(\bm{x}|y)=\mathcal{N}(\bm{x};\bm{\mu}_y,\Sigma_y)
\end{align*}

\section*{Visual page}
Draw data, model, loss surface, gradient step, and validation error.

\section*{GoodNotes pages}
\begin{itemize}[leftmargin=1.4em]
  \item Page ST-ML001: empirical risk
  \item Page ST-ML002: gradient descent
  \item Page ST-ML003: probabilistic model
\end{itemize}

\section*{Web interaction}
A loss-landscape slider that changes learning rate and overfitting.

\section*{Obsidian and Notion}
\begin{tabular}{p{0.22\linewidth}p{0.68\linewidth}}
\toprule
Layer & Output \\
\midrule
Obsidian & Stanford ML -> CS229 \\
Notion & Problem-set status and formula mastery. \\
\bottomrule
\end{tabular}

\section*{First study move}
Open the official source, copy the prerequisite checkpoint into GoodNotes, then write one formula in your own words before watching or reading more.

\end{document}
